% --- Archon physics formalization source begin ---
% archon:physics
% archon:covers PhyXMiniProblems/problem_phyx_mini_0652.lean
% archon:source-report reports/phyx_mini/problem_phyx_mini_0652.source.json

\chapter{Physics problem phyx\_mini\_0652}
\label{ch:PhyXMiniProblems_problem_phyx_mini_0652}

\paragraph{Problem source.}
\#\# Physical scenario

Figure shows the wave function of a particle confined between \$x = -4.0\$ \$mm\$ and \$x = 4.0\$ \$mm\$. The wave function is zero outside this region.

\#\# Question

Calculate the probability of finding the particle in the interval \$-2.0\$ \$mm\$ \$\textbackslash{}le x \textbackslash{}le 2.0\$ \$mm\$.

\#\# Answer choices

- A: A: 0.08

- B: B: 0.1

- C: C: 0.125

- D: D: 0.25

\#\# Figure caption (auxiliary; use the image as primary evidence)

The image is a plot of a piecewise linear function \textbackslash{}(\textbackslash{}psi(x)\textbackslash{}) against the \textbackslash{}(x\textbackslash{})-axis, measured in millimeters (mm).

\#\#\# Details:

1. **Axes:**
   - The horizontal axis is labeled \textbackslash{}(x \textbackslash{}, (\textbackslash{}text\{mm\})\textbackslash{}).
   - The vertical axis is labeled \textbackslash{}(\textbackslash{}psi(x)\textbackslash{}).

2. **Ticks and Labels:**
   - The \textbackslash{}(x\textbackslash{})-axis has ticks at -4, -2, 0, 2, and 4.
   - The \textbackslash{}(y\textbackslash{})-axis has labels at \textbackslash{}(c\textbackslash{}) and \textbackslash{}(-c\textbackslash{}).

3. **Function Behavior:**
   - The function is a piecewise linear line drawn in red.
   - It starts at the point (-4, -c).
   - It connects to the point (0, 0).
   - The function rises linearly to (4, c).
   - Beyond \textbackslash{}(x = 4\textbackslash{}) on the graph, the line drops vertically to touch \textbackslash{}(x = 4\textbackslash{}).

4. **Function Shape:**
   - The function forms a triangular shape with a peak at (4, c).
   - On the left, it forms an open vertical segment at \textbackslash{}(x = -4\textbackslash{}) from \textbackslash{}(-c\textbackslash{}) to a small finite value slightly below the \textbackslash{}(x\textbackslash{})-axis.

\#\#\# Observations:
- The graph appears to represent a symmetric triangular waveform centered at 0 with a span from \textbackslash{}(x = -4\textbackslash{}) to 4.

\#\# Dataset metadata

Quantum Phenomena; Numerical Reasoning, Physical Model Grounding Reasoning

\paragraph{Recorded answer/context.}
C: C: 0.125

\paragraph{Figure/image path.}
/root/proposal\_for\_physic/science-mango/phyx\_mini\_run/phyx\_data/test\_image/652.png

\paragraph{Formalization target.}
create a compiling Lean file with sorry bodies at `PhyXMiniProblems/problem\_phyx\_mini\_0652.lean`. The Lean declarations must preserve the physical quantities, dimensions or dimensional roles, figure labels, governing-law hypotheses, and final relation expressed by this problem.
Use Mathlib/Physlib names found through LeanExplore where available. If a physics API is missing, introduce faithful local abstractions rather than scalar placeholder aliases.

\begin{theorem}[Physics formalization target]
\label{thm:physics:phyx_mini_0652:target}
\lean{PhyXMiniProblems.ProblemPhyXMini0652.problem_phyx_mini_0652}
\uses{def:physics:phyx-mini-0652:phyxminiproblems-problemphyxmini0652-confinedparticlepositionsetup, def:physics:phyx-mini-0652:phyxminiproblems-problemphyxmini0652-requestedpositionprobability, def:physics:phyx-mini-0652:phyxminiproblems-problemphyxmini0652-matchesconfinedparticlescenario, def:physics:phyx-mini-0652:phyxminiproblems-problemphyxmini0652-matchessuppliedwavefunctionfigure, def:physics:phyx-mini-0652:phyxminiproblems-problemphyxmini0652-hasphysicalwavefunctionparameters, def:physics:phyx-mini-0652:phyxminiproblems-problemphyxmini0652-satisfiesnormalizedbornpositionlaw, def:physics:phyx-mini-0652:phyxminiproblems-problemphyxmini0652-recordeddatasetanswer, def:physics:phyx-mini-0652:phyxminiproblems-problemphyxmini0652-isuniquematchingdisplayedprobability}
The assigned autoformalize agent should translate this physics problem into Lean declarations in the covered file, with theorem and lemma proof bodies written as `by sorry`.
\end{theorem}
\begin{proof}
This is an autoformalization task, not a proof task. Produce statements that can later be proved by the physics prover without weakening the physical model.
\end{proof}

% --- Archon declaration topology begin ---
\section{Declaration topology}
\begin{definition}[Wave Function Figure Label]
  \label{def:physics:phyx-mini-0652:phyxminiproblems-problemphyxmini0652-wavefunctionfigurelabel}
  \lean{PhyXMiniProblems.ProblemPhyXMini0652.WaveFunctionFigureLabel}
  Labels printed on the axes and vertical scale of the supplied figure
\end{definition}
\begin{proof}
  This is the declared model definition of Wave Function Figure Label.
\end{proof}
\begin{definition}[Horizontal Tick]
  \label{def:physics:phyx-mini-0652:phyxminiproblems-problemphyxmini0652-horizontaltick}
  \lean{PhyXMiniProblems.ProblemPhyXMini0652.HorizontalTick}
  The five labelled horizontal-axis ticks visible in the figure
\end{definition}
\begin{proof}
  This is the declared model definition of Horizontal Tick.
\end{proof}
\begin{definition}[Supplied Wave Function Figure]
  \label{def:physics:phyx-mini-0652:phyxminiproblems-problemphyxmini0652-suppliedwavefunctionfigure}
  \lean{PhyXMiniProblems.ProblemPhyXMini0652.SuppliedWaveFunctionFigure}
  \uses{def:physics:phyx-mini-0652:phyxminiproblems-problemphyxmini0652-wavefunctionfigurelabel, def:physics:phyx-mini-0652:phyxminiproblems-problemphyxmini0652-horizontaltick}
  Raw quantitative and qualitative evidence transcribed from 652.png. The plotted amplitude is a real readout of the complex wavefunction. This structure stores no interval probability and no answer-choice value
\end{definition}
\begin{proof}
  This is the declared model definition of Supplied Wave Function Figure.
\end{proof}
\begin{definition}[Confined Particle Position Setup]
  \label{def:physics:phyx-mini-0652:phyxminiproblems-problemphyxmini0652-confinedparticlepositionsetup}
  \lean{PhyXMiniProblems.ProblemPhyXMini0652.ConfinedParticlePositionSetup}
  \uses{def:physics:phyx-mini-0652:phyxminiproblems-problemphyxmini0652-suppliedwavefunctionfigure}
  A particle confined to a one-dimensional interval, together with the representative shown in the figure and its position probability observable. All fields ending in Millimeters or PerSqrtMillimeter are calibrated scalar readouts in the units printed or implied by the graph. In particular, the wavefunction itself is not replaced by a scalar physical quantity: it is a complex function with the appropriate amplitude role
\end{definition}
\begin{proof}
  This is the declared model definition of Confined Particle Position Setup.
\end{proof}
\begin{definition}[position Density Per Millimeter]
  \label{def:physics:phyx-mini-0652:phyxminiproblems-problemphyxmini0652-positiondensitypermillimeter}
  \lean{PhyXMiniProblems.ProblemPhyXMini0652.positionDensityPerMillimeter}
  \uses{def:physics:phyx-mini-0652:phyxminiproblems-problemphyxmini0652-confinedparticlepositionsetup}
  The Born position-density readout in inverse millimetres
\end{definition}
\begin{proof}
  This is the declared model definition of position Density Per Millimeter.
\end{proof}
\begin{definition}[requested Position Region]
  \label{def:physics:phyx-mini-0652:phyxminiproblems-problemphyxmini0652-requestedpositionregion}
  \lean{PhyXMiniProblems.ProblemPhyXMini0652.requestedPositionRegion}
  \uses{def:physics:phyx-mini-0652:phyxminiproblems-problemphyxmini0652-confinedparticlepositionsetup}
  The closed position interval whose probability is requested
\end{definition}
\begin{proof}
  This is the declared model definition of requested Position Region.
\end{proof}
\begin{definition}[requested Position Probability]
  \label{def:physics:phyx-mini-0652:phyxminiproblems-problemphyxmini0652-requestedpositionprobability}
  \lean{PhyXMiniProblems.ProblemPhyXMini0652.requestedPositionProbability}
  \uses{def:physics:phyx-mini-0652:phyxminiproblems-problemphyxmini0652-confinedparticlepositionsetup, def:physics:phyx-mini-0652:phyxminiproblems-problemphyxmini0652-requestedpositionregion}
  The dimensionless probability of the requested position event
\end{definition}
\begin{proof}
  This is the declared model definition of requested Position Probability.
\end{proof}
\begin{definition}[Matches Confined Particle Scenario]
  \label{def:physics:phyx-mini-0652:phyxminiproblems-problemphyxmini0652-matchesconfinedparticlescenario}
  \lean{PhyXMiniProblems.ProblemPhyXMini0652.MatchesConfinedParticleScenario}
  \uses{def:physics:phyx-mini-0652:phyxminiproblems-problemphyxmini0652-confinedparticlepositionsetup}
  Numerical data stated in the problem text, including confinement to [-4 mm, 4 mm], the query interval [-2 mm, 2 mm], and vanishing amplitude outside the confinement region. No probability value occurs here
\end{definition}
\begin{proof}
  This is the declared model definition of Matches Confined Particle Scenario.
\end{proof}
\begin{definition}[Matches Supplied Wave Function Figure]
  \label{def:physics:phyx-mini-0652:phyxminiproblems-problemphyxmini0652-matchessuppliedwavefunctionfigure}
  \lean{PhyXMiniProblems.ProblemPhyXMini0652.MatchesSuppliedWaveFunctionFigure}
  \uses{def:physics:phyx-mini-0652:phyxminiproblems-problemphyxmini0652-confinedparticlepositionsetup}
  Exact labels, tick locations, and line profile read from the supplied raster. The interior profile is the straight line ψ(x) = c x / (4 mm) in millimetre-coordinate readouts. It is a figure/model premise and contains neither the integral over the queried interval nor its numerical value
\end{definition}
\begin{proof}
  This is the declared model definition of Matches Supplied Wave Function Figure.
\end{proof}
\begin{definition}[Has Physical Wave Function Parameters]
  \label{def:physics:phyx-mini-0652:phyxminiproblems-problemphyxmini0652-hasphysicalwavefunctionparameters}
  \lean{PhyXMiniProblems.ProblemPhyXMini0652.HasPhysicalWaveFunctionParameters}
  \uses{def:physics:phyx-mini-0652:phyxminiproblems-problemphyxmini0652-confinedparticlepositionsetup}
  Positivity and ordering conditions implicit in the physical setup
\end{definition}
\begin{proof}
  This is the declared model definition of Has Physical Wave Function Parameters.
\end{proof}
\begin{definition}[Satisfies Normalized Born Position Law]
  \label{def:physics:phyx-mini-0652:phyxminiproblems-problemphyxmini0652-satisfiesnormalizedbornpositionlaw}
  \lean{PhyXMiniProblems.ProblemPhyXMini0652.SatisfiesNormalizedBornPositionLaw}
  \uses{def:physics:phyx-mini-0652:phyxminiproblems-problemphyxmini0652-confinedparticlepositionsetup, def:physics:phyx-mini-0652:phyxminiproblems-problemphyxmini0652-positiondensitypermillimeter}
  The normalized one-dimensional Born rule. MemHS is Physlib's membership predicate for the one-dimensional L² Hilbert space. The probability equation is stated for every measurable event, so it is a governing law rather than the numerical result for the particular interval in this problem
\end{definition}
\begin{proof}
  This is the declared model definition of Satisfies Normalized Born Position Law.
\end{proof}
\begin{definition}[Answer Choice]
  \label{def:physics:phyx-mini-0652:phyxminiproblems-problemphyxmini0652-answerchoice}
  \lean{PhyXMiniProblems.ProblemPhyXMini0652.AnswerChoice}
  Labels of the four probability choices printed with the problem
\end{definition}
\begin{proof}
  This is the declared model definition of Answer Choice.
\end{proof}
\begin{definition}[displayed Probability]
  \label{def:physics:phyx-mini-0652:phyxminiproblems-problemphyxmini0652-displayedprobability}
  \lean{PhyXMiniProblems.ProblemPhyXMini0652.displayedProbability}
  \uses{def:physics:phyx-mini-0652:phyxminiproblems-problemphyxmini0652-answerchoice}
  The dimensionless probability displayed beside each answer label
\end{definition}
\begin{proof}
  This is the declared model definition of displayed Probability.
\end{proof}
\begin{definition}[recorded Dataset Answer]
  \label{def:physics:phyx-mini-0652:phyxminiproblems-problemphyxmini0652-recordeddatasetanswer}
  \lean{PhyXMiniProblems.ProblemPhyXMini0652.recordedDatasetAnswer}
  \uses{def:physics:phyx-mini-0652:phyxminiproblems-problemphyxmini0652-answerchoice}
  The answer label recorded by the supplied dataset; it is never a premise
\end{definition}
\begin{proof}
  This is the declared model definition of recorded Dataset Answer.
\end{proof}
\begin{definition}[Matches Displayed Probability]
  \label{def:physics:phyx-mini-0652:phyxminiproblems-problemphyxmini0652-matchesdisplayedprobability}
  \lean{PhyXMiniProblems.ProblemPhyXMini0652.MatchesDisplayedProbability}
  \uses{def:physics:phyx-mini-0652:phyxminiproblems-problemphyxmini0652-confinedparticlepositionsetup, def:physics:phyx-mini-0652:phyxminiproblems-problemphyxmini0652-requestedpositionprobability, def:physics:phyx-mini-0652:phyxminiproblems-problemphyxmini0652-answerchoice, def:physics:phyx-mini-0652:phyxminiproblems-problemphyxmini0652-displayedprobability}
  A choice matches when its displayed value is the modeled event probability
\end{definition}
\begin{proof}
  This is the declared model definition of Matches Displayed Probability.
\end{proof}
\begin{definition}[Is Unique Matching Displayed Probability]
  \label{def:physics:phyx-mini-0652:phyxminiproblems-problemphyxmini0652-isuniquematchingdisplayedprobability}
  \lean{PhyXMiniProblems.ProblemPhyXMini0652.IsUniqueMatchingDisplayedProbability}
  \uses{def:physics:phyx-mini-0652:phyxminiproblems-problemphyxmini0652-confinedparticlepositionsetup, def:physics:phyx-mini-0652:phyxminiproblems-problemphyxmini0652-answerchoice, def:physics:phyx-mini-0652:phyxminiproblems-problemphyxmini0652-matchesdisplayedprobability}
  A choice is the unique displayed value matching the modeled probability
\end{definition}
\begin{proof}
  This is the declared model definition of Is Unique Matching Displayed Probability.
\end{proof}
\begin{lemma}[central Interval Probability eq one Eighth]
  \label{lem:physics:phyx-mini-0652:phyxminiproblems-problemphyxmini0652-centralintervalprobability-eq-oneeighth}
  \lean{PhyXMiniProblems.ProblemPhyXMini0652.centralIntervalProbability_eq_oneEighth}
  \uses{def:physics:phyx-mini-0652:phyxminiproblems-problemphyxmini0652-confinedparticlepositionsetup, def:physics:phyx-mini-0652:phyxminiproblems-problemphyxmini0652-requestedpositionprobability, def:physics:phyx-mini-0652:phyxminiproblems-problemphyxmini0652-matchesconfinedparticlescenario, def:physics:phyx-mini-0652:phyxminiproblems-problemphyxmini0652-matchessuppliedwavefunctionfigure, def:physics:phyx-mini-0652:phyxminiproblems-problemphyxmini0652-hasphysicalwavefunctionparameters, def:physics:phyx-mini-0652:phyxminiproblems-problemphyxmini0652-satisfiesnormalizedbornpositionlaw}
  Integrating the squared linear wavefunction over [-2 mm, 2 mm] and using normalization over [-4 mm, 4 mm] gives the exact probability 1/8. This is a derived result, not a premise
\end{lemma}
\begin{proof}
  The conclusion follows from the governing relations and figure readouts named in the statement of central Interval Probability eq one Eighth.
\end{proof}
% --- Archon declaration topology end ---
% --- Archon physics formalization source end ---
